\documentclass[11pt]{article}
\usepackage[ngerman]{babel}
\usepackage[a4paper, left=2.5cm, right=2.5cm, top=3.5cm]{geometry}
\usepackage{amsmath}
\usepackage{nccmath}
\usepackage{amssymb}
\usepackage[pdftex]{graphicx}
\usepackage{listings}
\usepackage{booktabs}
\usepackage{tikz}
\usepackage{enumitem}
\usepackage{fancyhdr}
\usepackage{geometry}
\usepackage{mathtools}
\usepackage{lipsum}
\usepackage{amsthm}
\usepackage{amsfonts}
\usepackage{geometry}
\usepackage[dvipsnames]{xcolor}
\usepackage{enumitem}
\usepackage{mathtools}
\usepackage{fancyhdr}
\usepackage{listings}
\usepackage{tikz}
\usepackage[utf8]{inputenc}
\newcommand{\bbN}{\mathbbm{N}}
\newcommand{\bbZ}{\mathbbm{Z}}
\newcommand{\bbQ}{\mathbbm{Q}}
\newcommand{\bbR}{\mathbbm{R}}
\newcommand{\bbC}{\mathbbm{C}}
\newcommand{\bbI}{\mathbbm{I}}
\newcommand{\bigO}{\mathcal{O}}
\newcommand{\bigT}{\mathcal{T}}
\newcommand\numberthis{\addtocounter{equation}{1}\tag{\theequation}}
\newcommand\abv[2]{\stackrel{\mathclap{\normalfont\mbox\tiny{#2}}}{#1}}
\newcommand\norm[1]{\left\lVert{#1}\right\rVert}
\lstset{basicstyle=\ttfamily, mathescape}
\geometry{top=2.5cm, bottom=2.5cm, left=2.5cm, right=2.5cm}
\setlength{\parindent}{0 pt}
\hbadness=99999 % No fucking underfull hbox warning
\hfuzz=9999pt % No fucking overfull hbox warning
\pagestyle{fancy}
\tikzstyle{textbox} = [draw, rounded corners, minimum height=2em, minimum width=4em]
\tikzstyle{arrow} = [->, thick]
\setlength{\headheight}{42.91258pt}
\lstdefinestyle{pretty}{
  backgroundcolor=\color{gray!10},
  frame=single,
  rulecolor=\color{gray!60},
  basicstyle=\ttfamily\small,
  keywordstyle=\color{blue}\bfseries,
  commentstyle=\color{green!60!black}\itshape,
  stringstyle=\color{red!70!black},
  numberstyle=\tiny\color{gray},
  numbers=left,
  stepnumber=1,
  numbersep=8pt,
  tabsize=2,
  breaklines=true,
  breakatwhitespace=true,
  showstringspaces=false,
  captionpos=b
}

\lstset{style=pretty}
\lstset{escapeinside={(*@}{@*)}}

\lhead{Computergraphik\\
Wintersemester 2025/26\\
\today}
\chead{\huge\textbf{Übungsblatt 3}}
\rhead{Maximilian Peresunchak 3232875\\
Nico Reng 3731402\\
Viorel Tsigos 3720183}
\begin{document}
\section*{Aufgabe 1:}
\begin{enumerate}
  \item Folgende Sachen haben wir gegeben: \begin{align*}
  &\text{Ortsvektor des Ausgangspunkts: } p = (-1,~-1,~1)^T \\
  &\text{Richtungsvektor: } d = (1,~1,~0)^T \\
  &\text{Kugel: } r = 3, ~\text{Mittelpunkt: } M = (5,~5,~2)^T \\
  &r(t) = p~+~td
\end{align*}
Wir wissen auch dass die Kugelgleichung $(r-M) \cdot (r-M)$ lautet. Wir setzen hier jetzt $r(t)$ in die Gleichung ein und definieren uns einen Hilfsvektor $v=p-M$:
\begin{align*}
  ((p + td) - M) \cdot ((p + td) - M) &= r^2 \\
  t^2(d \cdot d) + 2t(d \cdot v) + (v \cdot v) &= r^2 \\
  t^2(d \cdot d) + 2t(d \cdot v) + (v \cdot v - r^2) &= 0 
\end{align*}
Jetzt müssen wir jetzt noch den Hilfsvektor $v$ berechnen: 
\begin{align*}
  v = p-m = \begin{pmatrix*}
    -1 \\
    -1 \\
    1   
  \end{pmatrix*} - \begin{pmatrix*}
    5\\
    5\\ 
    2
  \end{pmatrix*} = \begin{pmatrix*}
    -6\\
    -6\\
    -1
  \end{pmatrix*}
\end{align*}
Daraus folgen nun die Koeffizieten der quadratischen Gleichung:
\begin{align*}
  A &= (d \cdot d) = \begin{pmatrix*}
    1\\
    1\\
    0
  \end{pmatrix*} \cdot \begin{pmatrix*}
    1\\
    1\\
    0
  \end{pmatrix*} = 1^2 + 1^2 + 0^2 = 2\\\\
  B &= 2(d \cdot v) = 2 \cdot \left( \begin{pmatrix*}
    1\\
    1\\
    0
  \end{pmatrix*} \cdot \begin{pmatrix*}
    -6\\
    -6\\
    -1
  \end{pmatrix*} \right) = 2 \cdot (-6 -6 + 0) = -24\\ \\
  C &= (v \cdot v - r^2) = \left( \begin{pmatrix*}
      -6\\
      -6\\
      -1
  \end{pmatrix*} \cdot \begin{pmatrix*}
     -6\\
    -6\\
    -1
  \end{pmatrix*} - 3^2 \right) 
  \\ &= (-6)^2 + (-6)^2 + (-1)^2 - 3^2 = 36 + 36 + 1 - 9 = 64
\end{align*}
Daraus folgt nun die quadratische Gleichung, die wir danach lösen: 
\begin{align*}
  2t^2 - 24t + 64 = 0
   \Rightarrow t_{1,2} &= \frac{-b \pm \sqrt{b^2 - 4ac}}{2a}\\ &= \frac{24 \pm \sqrt{(-24)^2 - 4 \cdot 2 \cdot 64}}{2 \cdot 2}\\&= \frac{24 \pm \sqrt{576 - 512}}{4} \\&= \frac{24 \pm \sqrt{64}}{4}\\ &= \frac{24 \pm 8}{4} \\
   &\Rightarrow t_1 = 8, ~t_2 = 4
\end{align*}
Das setzen wir jetzt in $r(t) = p + td$ ein um die Schnittpunkte zu berechnen: 
\begin{align*}
  r(4) &= \begin{pmatrix*}
    -1\\
    -1\\
    1
  \end{pmatrix*} + 4 \cdot \begin{pmatrix*}
    1\\
    1\\
    0
  \end{pmatrix*} = \begin{pmatrix*}
    3\\
    3\\
    1
  \end{pmatrix*} \\\\
  r(8) &= \begin{pmatrix*}
    -1\\
    -1\\
    1
  \end{pmatrix*} + 8 \cdot \begin{pmatrix*}
    1\\
    1\\
    0
  \end{pmatrix*} = \begin{pmatrix*}
    7\\
    7\\
    1
  \end{pmatrix*} \\
\end{align*} 
Das sind jetzt die beiden Schnittpunkte. Hier sieht man jetzt relativ schnell, welches der Schnittpunkt für den Betrachter sichtbar ist. Nämlich der Schnittpunkt mit dem kleinsten \textbf{positiven} Parameterwert. Hier ist es einmal $t_1 = 8$ und $t_2 = 4$. Somit ist der sichtbare Schnittpunkt $r(4) = (3,~3,~1)^T$.
\item Neu gegeben jetzt: $d = 3$ also Abstand vom Ursprung und den Normalenvektor $n$ in Richtung $(1, 1, 0)^T$
Wir normieren jetzt erst mal unseren Normalenvektor $n$: 
\begin{align*}
  &|n| = \sqrt{1^2 + 1^2 + 0^2} = \sqrt{2} \\
  &\text{Kehrwert: } \frac{1}{\sqrt{2}} 
\end{align*}
Daraus folgt jetzt der normierte Normalenvektor:
\begin{align*} \frac{1}{\sqrt{2}} \cdot 
  \begin{pmatrix*}
    1 \\
    1\\
    0
  \end{pmatrix*}
\end{align*}
Wir wissen auch dass die Ebenengleichung $x \cdot n_0 - d = 0$ ist. Wir setzen jetzt für die Schnittberechnung den Strahl $r(t) = p + td$ in die Ebenengleichung ein:
\begin{align*}
  (p + td) \cdot n_0 - d &= 0 \\
  (p \cdot n_0) + t(d \cdot n_0) - d &= 0 \\
  t(d \cdot n_0) &= d - (p \cdot n_0) \\
  t &= \frac{d - (p \cdot n_0)}{(d \cdot n_0)}
\end{align*}
Jetzt bestimmen wir die Skalarprodukte: 
\begin{align*}
  1.~~p \cdot n_0 &= \begin{pmatrix*}
    -1\\
    -1\\
    1
  \end{pmatrix*} \cdot \left( \frac{1}{\sqrt{2}} \cdot \begin{pmatrix*}
    1\\
    1\\
    0
  \end{pmatrix*} \right) = \frac{1}{\sqrt{2}} \cdot (-1 -1 + 0) = \frac{-2}{\sqrt{2}} = -\sqrt{2} \\ \\
  2.~~d \cdot n_0 &= \begin{pmatrix*}
    1\\
    1\\
    0
  \end{pmatrix*} \cdot \left( \frac{1}{\sqrt{2}} \cdot \begin{pmatrix*}
    1\\
    1\\
    0
  \end{pmatrix*} \right) = \frac{1}{\sqrt{2}} \cdot (1 + 1 + 0) = \frac{2}{\sqrt{2}} = \sqrt{2}
\end{align*}
Jetzt können wir den Schnittpunkt berechnen: 
\begin{align*}
  r(t) &= \begin{pmatrix*}
    -1\\
    -1\\
    1
  \end{pmatrix*} + \left( \frac{3}{\sqrt{2}} + 1 \right) \cdot \begin{pmatrix*}
    1\\
    1\\
    0
  \end{pmatrix*} = \begin{pmatrix*}
    -1 +\frac{3}{\sqrt{2}} + 1\\
    -1 +\frac{3}{\sqrt{2}} + 1\\
    1 + 0
  \end{pmatrix*}  = \begin{pmatrix*}
    \frac{3}{\sqrt{2}} \\
    \frac{3}{\sqrt{2}} \\
    1
  \end{pmatrix*}
\end{align*} \\
\item Gegeben ist ein Dreieck mit drei Eckpunkten: $A = (6, 0, 0) ^T$, $B = (0, 6, 0)^T$, $C = (0, 0, 6)^T$
Die drei Punkte spannen also eine Ebene auf. Daraus folgt dann folgende Gleichung: \\$x+y+z = 6$. Wir können hier relativ einfach den Normalvektor $n_E$ ablesen: \begin{align*}
    n_E = \begin{pmatrix*}
      1\\
      1\\
      1
    \end{pmatrix*}
\end{align*}
Wir wissen auch die Ebenengleichung: $x \cdot n_E = 6$. \\
Wir setzen jetzt für $x$ den Punkt $r(t)$ ein und stellen nach $t$ um: 
\begin{align*}
  (p + td) \cdot n_E &= 6 \\
  (p \cdot n_E) + t(d \cdot n_E) &= 6 \\
  t(d \cdot n_E) &= 6 - (p \cdot n_E) \\
  t &= \frac{6 - (p \cdot n_E)}{(d \cdot n_E)}
\end{align*}
Das ist solange gültig, solange $(d \cdot n_E) \neq 0$ ist, was sozusagen heißt dass der Strahl nicht parallel zur Ebene läuft. Jetzt rechnen wir wieder die Skalarprodukte aus:
\begin{align*}
  1.~~p \cdot n_E &= \begin{pmatrix*}
    -1\\
    -1\\
    1
  \end{pmatrix*} \cdot \begin{pmatrix*}
    1\\
    1\\
    1
  \end{pmatrix*} = -1 \cdot 1 + -1 \cdot 1 + 1 \cdot 1 = -1 \\ \\
  2.~~d \cdot n_E &= \begin{pmatrix*}
    1\\
    1\\
    0
  \end{pmatrix*} \cdot \begin{pmatrix*}
    1\\
    1\\
    1
  \end{pmatrix*} = 1 \cdot 1 + 1 \cdot 1 + 1 \cdot 0 = 2
\end{align*}
Jetzt können wir $t$ bestimmen: 
\begin{align*}
  t &= \frac{6 - (-1)}{2} = \frac{7}{2} = 3.5
\end{align*}
Jetzt setzen wir das $t$ in die Strahlengleichung $r(t)$ ein um den Schnittpunkt zu bestimmen:
\begin{align*}
  r(3.5) &= \begin{pmatrix*}
    -1\\
    -1\\
    1
  \end{pmatrix*} + 3.5 \cdot \begin{pmatrix*}
    1\\
    1\\
    0
  \end{pmatrix*} = \begin{pmatrix*}
    -1 + 3.5\\
    -1 + 3.5\\
    1 + 0
  \end{pmatrix*} = \begin{pmatrix*}
    2.5\\
    2.5\\
    1
  \end{pmatrix*}
\end{align*} \\ \\
Also ist der Schnittpunkt $S = \begin{pmatrix*}
    2.5\\
    2.5\\
    1
  \end{pmatrix*}$.
\item Dafür können wir die Baryzentrischen Koordinaten benutzen $S = \alpha A + \beta B + \delta C$. Folgende zwei Sachen müssen dafür gelten, dass der Sichtstrahl das Dreieck schneidet:
\begin{align*}
  &1.~~\alpha, \beta, \gamma \geq 0 \\
  &2.~~\alpha + \beta + \gamma = 1
\end{align*}
Setzen wir das jetzt ein, kommt folgendes raus:
\begin{align*}
 \begin{pmatrix*}
    2,5 \\
    2,5 \\
    1
  \end{pmatrix*} &= \alpha \cdot \begin{pmatrix*}
    6 \\
    0 \\
    0
  \end{pmatrix*} + \beta \cdot \begin{pmatrix*}
    0 \\
    6 \\
    0
  \end{pmatrix*} + \gamma \cdot \begin{pmatrix*}
    0 \\
    0 \\
    6
 \end{pmatrix*}
\end{align*}
Daraus folgen jetzt drei Gleichungen die wir auflösen: 
\begin{align*}
  &1.~~6\alpha =  2.5 \Rightarrow \alpha = \frac{2.5}{6} \\
  &2.~~6\beta  = 2.5 \Rightarrow \beta = \frac{2.5}{6} \\
  &3.~~6\gamma =  1 \Rightarrow \gamma = \frac{1}{6}
\end{align*}
Jetzt überprüfen wir noch die beiden Bedingungen:
\begin{align*}  
  &1.~~\alpha, \beta, \gamma \geq 0~~~~ \checkmark \\
  &2.~~\alpha + \beta + \gamma = \frac{2.5}{6} + \frac{2.5}{6} + \frac{1}{6} = \frac{6}{6} = 1~~~~ \checkmark
\end{align*}
Somit schneidet der Sichtstrahl das Dreieck.
\end{enumerate}
\end{document}